\section{Tracing Frameworks}

\subsection{Ftrace}

Structure:

- A bit of history
- architecture
- in-kernel api
- user space api

Good intro:
\url{https://docs.redhat.com/en/documentation/red_hat_enterprise_linux_for_real_time/7/html/tuning_guide/appe-detailed_description_of_ftrace}

There is a naming issue as ftrace was born as the function tracer, but has
evolved into the main tracing infrastructure used by other types of tracers. We
will now talk about the general unfrastructure.

We write to buffer with \verb|trace_printk()|.  SEE:
\verb|kernel/trace/trace_printk.c| modules \verb|trace_printk()|'s formats are
autosaved in \verb|struct trace_bprintk_fmt| which are queued on
\verb|trace_bprintk_fmt_list|.

An important abstraction is the struct tracer (kernel/trace/trace.h) which
represents a specific tracer and its callbacks to interact with tracefs. Then
you do \verb|tracer_init| and other \verb|tracer_*| functions.

trace.h:1450

   \verb|extern struct tracer nop_trace;|

From: \url{https://www.kernel.org/doc/html/latest/trace/events.html} Trace
events can be made to conditionally invoke trigger ‘commands’ which can take
various forms and are described in detail below; examples would be enabling or
disabling other trace events or invoking a stack trace whenever the trace event
is hit. Whenever a trace event with attached triggers is invoked, the set of
trigger commands associated with that event is invoked.

\begin{verbatim}

trace.h:1997
   struct event_command {

\end{verbatim}

Not that you also have event filtering.

The trace array contains per-CPU trace arrays. This is the highest level data
structure that indivisual tracers deal with.

Interesting: lockless ringbuffer design TODO:
\url{https://www.kernel.org/doc/html/latest/trace/ring-buffer-design.html} The
ring buffer is divided in pages. Each CPU has its own list of pages where it can
write, and the reader may read from any per-cpu buffer. For this reason for each
per-cpu buffer, the reader has its own page separate from the rest. When the
reader has done reading it, it swaps it with another page. If the writer wraps,
it needs to push the head forward. Here we need to use cmpxchg since we need to
make sure that the writer and the reader are not changing the head at the same
time.  SEE: \verb|kernel/trace/ring_buffer.c|

There are several tracers here, the most famous one is the function
tracer.

\subsection{perf\_event}

Structure:

- A bit of history
- architecture
- in-kernel api
- user space api

Perf: event state machine

You interact with perf via the \verb|perf_event_open| syscall.
Used for performance counting. Has its own ring buffer.

They are simple as they only signal that an event is happened. This
is then processed by another program like perf.

\subsection{eBPF}

Structure:

- A bit of history
- architecture
- in-kernel api
- user space api

TODO: just an introduction.
Has its own ring buffer, but also provides maps which are of
various types.

\subsection{audit}

TODO

\subsection{Vendor Specific Systems}

TODO

\subsection{Out of Tree}

  - systemtap
  - LTTng
  - ktap
  - dtrace4linux
  - DTrace
  - sysdig
  - ...
